\section{Decreasing limits under the envelope property}
\label{sec:ep-decreasing-limits}

This section supplies the arbitrary-class decreasing-limit step separately from
the homogeneous decreasing-limit theorem. The envelope property is used only
at the limiting class. On unibranch models it is supplied by the main theorem,
with the same mathematical inputs as that theorem.

\begin{lemma}[Convex rays of decreasing limits]\label{lem:ep-limit-convex-rays}
\lean{EnvelopeProperty.MultiModel.iInf_psh_ne_top, EnvelopeProperty.MultiModel.iInf_psh_le_triv, EnvelopeProperty.MultiModel.iInf_psh_triv_ne_bot, EnvelopeProperty.MultiModel.scaleFun_iInf_psh_convex, EnvelopeProperty.MultiModel.scaleFun_iInf_psh_ne_bot, EnvelopeProperty.MultiModel.convexOn_rayFun_iInf_psh}
\leanok
\uses{lem:gl-action-convex, lem:gl-scaling, thm:gl-fs-system-theory}
\source{fhx-envelope-20261028:prop:BJ-package, bj22-trivval:lem:actionconvex}
Let $u_i\in\PSH(\mathcal X,\theta_i)$ decrease to $u\not\equiv-\infty$.
Then $u\le u(\star)\in\R$. For every $v$, the function
$g_v(t)=t u(t^{-1}v)$ on $t>0$ satisfies the convexity inequality and
$g_v(t)\le t u(\star)$. It is either identically $-\infty$ or finite everywhere.
\end{lemma}

\begin{proof}
The inequalities $u_i(v)\le u_i(\star)$ pass to infima. A point where $u$ is
finite gives a finite lower bound for $u(\star)$; any one $u_i(\star)$ is a
finite upper bound. Scaling by a positive real commutes with arbitrary
infima. Thus $g_v$ is the decreasing limit of the convex rays of $u_i$.
The convexity inequality passes to this limit: both summands avoid $+\infty$,
so addition is continuous at their limiting values. If $g_v(t_0)>-\infty$,
write $t_0$ as a convex combination of $t$ and $t_0^2/t$, with coefficients
$t_0/(t+t_0)$ and $t/(t+t_0)$. A value $g_v(t)=-\infty$ would force
$g_v(t_0)=-\infty$. Hence all ray values are finite.
\end{proof}

\begin{lemma}[Reconstruction of a decreasing limit]\label{lem:ep-limit-reconstruction}
\lean{EnvelopeProperty.MultiModel.iInf_psh_eq_iSup_legendreSlice}
\leanok
\uses{lem:ep-limit-convex-rays, lem:gl-scaling, thm:gl-homog-decomposition}
\source{fhx-envelope-20261028:eq:Legendre-reconstruction, bj22-trivval:thm:homogpsh}
For the limit $u$ above, writing $b=u(\star)$ and
$\hat u_\lambda=\inf_{t>0}t(u(t^{-1}\cdot)-\lambda)$, one has
$u=\sup_{\lambda<b}(\hat u_\lambda+\lambda)$ pointwise.
\end{lemma}

\begin{proof}
Evaluation at $t=1$ gives the upper bound. Fix $v$ with $u(v)>-\infty$.
The preceding lemma makes $g(t)=t u(t^{-1}v)$ a finite convex function.
Choose a supporting slope $c$ at $1$, so
$g(t)\ge g(1)+c(t-1)$. The upper bound $g(t)\le tb$ forces $c\le b$:
if $c>b$, choosing $t>\max(1,(c-g(1))/(c-b))$ contradicts the two bounds.
For every $\epsilon>0$, put $\lambda=c-\epsilon<b$. The supporting inequality
gives $g(t)-t\lambda\ge g(1)-c$ for all $t>0$, whence
$\hat u_\lambda(v)+\lambda\ge g(1)-\epsilon$. Letting $\epsilon$ decrease
to zero proves the lower bound. If $u(v)=-\infty$, the upper bound alone
gives equality.
\end{proof}

\begin{theorem}[Decreasing limits under the envelope property]\label{thm:ep-decreasing-limits}
\lean{EnvelopeProperty.MultiModel.legendreSlice_iInf, EnvelopeProperty.MultiModel.isPSH_of_antitone_of_hasEnvelopeProperty}
\leanok
\uses{lem:ep-limit-reconstruction, thm:nn-psh2-hom, cor:gl-pshhom, def:gl-psh}
\source{fhx-envelope-20261028:prop:BJ-package, bj22-trivval:thm:psh12}
Suppose $\mathrm{EP}_{\mathcal X}(\theta)$ holds. If
$u_i\in\PSH(\mathcal X,\theta_i)$ decrease to $u\not\equiv-\infty$ and
$\theta_i\to\theta$, then $u\in\PSH(\mathcal X,\theta)$.
\end{theorem}

\begin{proof}
For $\lambda<u(\star)$ every $(\hat u_i)_\lambda$ is homogeneous psh of
class $\theta_i$. Homogenization commutes with arbitrary infima, including
the real translate by $-\lambda$, so these functions decrease to
$\hat u_\lambda$. Homogeneous decreasing limits give
$\hat u_\lambda\in\PSH_{\mathrm{hom}}(\mathcal X,\theta)$.
The translated slices form a nonempty family of $\theta$-psh functions,
bounded above by the real number $u(\star)$. Their supremum is $u$ by the
reconstruction lemma. The function $u$ is upper semicontinuous as an infimum
of upper semicontinuous functions. Applying the envelope property therefore
gives $u=u^\star\in\PSH(\mathcal X,\theta)$.
\end{proof}

\begin{corollary}[Decreasing limits on unibranch models]\label{cor:ep-unibranch-decreasing-limits}
\lean{EnvelopeProperty.MultiModel.isPSH_of_antitone_of_isUnibranch}
\leanok
\uses{thm:ep-decreasing-limits, thm:nn-main}
\source{fhx-envelope-20261028:prop:BJ-package, fhx-envelope-20261028:thm:intro-main, fhx-envelope-20261028:cor:unibranch-decreasing-stability}
Under the hypotheses of the main theorem, the preceding decreasing-limit
property holds at every real class on every unibranch model. This is
\cite[Corollary~8.4]{FHX} on models.
\end{corollary}

\begin{proof}
The main theorem supplies the envelope property at the limiting class.
Apply the preceding theorem.
\end{proof}
